\chapter{Towards Benchmarking Neuro-Symbolic Description Logic Reasoners}
\label{ch:nesybench}

In the previous chapters, we introduced benchmarking methodologies for conventional Description Logic reasoners operating in both static and stream-based environments. These benchmarks target symbolic reasoning systems whose conclusions are derived through formally defined logical inference and whose correctness is assessed using traditional reasoning criteria such as soundness, completeness, scalability, and query performance.

At the same time, Description Logic reasoning is undergoing a significant evolution. Recent advances in neuro-symbolic Artificial Intelligence have led to reasoning systems that combine symbolic knowledge with neural learning, enabling reasoning over noisy, incomplete, and large-scale knowledge graphs. Rather than replacing conventional symbolic reasoners, these approaches address complementary reasoning scenarios where purely symbolic methods may be difficult to apply efficiently.

This evolution also introduces new benchmarking challenges. Existing benchmark frameworks, including OWL2Bench and GenACT, are designed for conventional symbolic reasoners and therefore cannot adequately evaluate properties unique to neuro-symbolic systems, such as approximate reasoning, learned representations, robustness to imperfect knowledge, or statistical generalization. Consequently, benchmarking neuro-symbolic Description Logic reasoners requires evaluation methodologies that differ fundamentally from those developed for classical symbolic reasoning. This chapter's account of the resulting desiderata and methodology is based on our own published analysis of this problem~\cite{singh2025nesybenchmarking}.

This chapter explores this emerging research direction, and it does so by starting from evidence rather than argument. Section~\ref{nesy} surveys the neuro-symbolic Description Logic landscape and the diversity of existing approaches, including a benchmarking effort we co-authored specifically on robustness to noise~\cite{loesch2024nsorn}, which is related to this chapter but is not itself a contribution of this thesis. Section~\ref{study} then takes two systems already documented in that survey and actually runs them, using this thesis's own ontology-generation tooling together with a real ontology from the literature as input. That exercise surfaces, first-hand, several of the exact difficulties the rest of the chapter argues are structural to the problem, rather than incidental to any one system. Sections~\ref{desiderata} and~\ref{benchmark} then state the resulting desiderata and a possible benchmark-design methodology, in each case pointing back to the concrete difficulty in Section~\ref{study} (or, where applicable, in~\cite{loesch2024nsorn}) that motivates it. In doing so, the chapter establishes a foundation for extending the benchmarking principles developed throughout this thesis to the next generation of Description Logic reasoning systems.

\section{Neuro-Symbolic Reasoning in the Description Logics Setting}
\label{nesy}

Neuro-symbolic Artificial Intelligence (NeSyAI)~\cite{neuro_symbolic,neurosymbolic_AI} aims to bridge the gap between symbolic logic and neural network-based learning by integrating structured reasoning with data-driven pattern recognition. In the context of reasoning, neuro-symbolic approaches seek to combine the interpretability and formal guarantees of symbolic systems with the robustness and scalability of neural models. However, unlike classical reasoning systems, these approaches typically relax strict guarantees such as soundness and completeness in favor of approximate reasoning and generalization.

Progress in this field faces significant challenges precisely because it is still emerging, in contrast to areas with extensive research and well-established benchmarks. Several models (Graph Neural Networks (GNN)~\cite{GNN}, Logic Tensor Networks~\cite{LTN}), methodologies (Inductive Logic Programming~\cite{ILP}), and adjacent ideas (explainable AI~\cite{XAI}, zero-shot learning~\cite{zero-shot}) all contribute to this field's diversity. Based on the reciprocal relationships between neural and symbolic components and how they benefit each other, Henry Kautz categorizes neuro-symbolic AI systems into six distinct types~\cite{henry}:

\begin{itemize}
\item \textbf{Symbolic Neuro Symbolic}: Input and output are represented symbolically, such as with words or sequences of words. These symbols are converted into vectors using methods like word2vec~\cite{word2vec} and then fed into a neural network for processing.
\item \textbf{Symbolic[Neuro]}: Symbolic solvers use neural models internally for some functions, as seen in systems like AlphaGo~\cite{alphago}.
\item \textbf{Neuro|Symbolic}: Neural and symbolic systems collaborate to enhance specific tasks, as in the Neuro-Symbolic Concept-Learner~\cite{concept_learner}.
\item \textbf{Neuro:Symbolic $\rightarrow$ Neuro}: Symbolic rules are taken as input and compiled during training, integrating symbolic knowledge into the structure of neural models, as demonstrated in Deep Learning For Symbolic Mathematics~\cite{symbolicmaths}.
\item \textbf{$Neuro_{Symbolic}$}: Symbolic rules are transformed into templates for structures within the neural network, as in Logic Tensor Networks~\cite{logictensornetworks}.
\item \textbf{Neuro[Symbolic]}: Symbolic reasoning is embedded inside a neural engine, such as Graph Neural Networks~\cite{GNN}.
\end{itemize}

Each category represents a distinct integration strategy, and this diversity of strategies is itself part of why systematic evaluation is hard: different methods and criteria are used to evaluate neuro-symbolic reasoning systems (Table~\ref{tab:nesy}), and this lack of a common evaluation basis is a significant barrier to progress.

In recent years, there have been significant advances in neuro-symbolic reasoners for description logics (DLs)~\cite{dlbook}, the formal underpinning of the Web Ontology Language (OWL 2)~\cite{owl2}. Most of this work focuses on classification and consistency checking~\cite{Sweb_foundation,sota2,sota1}; other reasoning tasks — instance retrieval, query rewriting, materialization, abduction, and explanation generation — remain relatively unexplored, despite varying considerably in intrinsic difficulty and offering a promising avenue for further work.

Rather than tackling classification and consistency directly, research in this area typically decomposes them into class subsumption, class membership, and satisfiability tasks. Common techniques include geometric embeddings~\cite{ELEm,Mondal,why_settle_for_one,faithful_embeddings} that map ontological relationships to geometric spaces, and approaches that emulate logical reasoning through machine learning~\cite{emulation,pointer_nw,noise-tolerant}. A comprehensive overview of the state-of-the-art neuro-symbolic reasoning landscape is given in~\cite{sota1,sota2}. Outside these categories, a smaller body of work merges neuro-symbolic reasoning with query rewriting~\cite{query_rewriting} — a Knowledge Graph (KG)~\cite{KG}-enhanced neural approach that integrates auxiliary knowledge from a product KG to improve query reformulation for e-commerce search.

Existing traditional benchmarks such as LUBM~\cite{LUBM}, UOBM~\cite{UOBM}, and OWL2Bench~\cite{owl2bench} are not well suited to evaluating neuro-symbolic reasoners, owing to their narrow focus on conventional reasoning tasks and metrics such as reasoning time that do not align with neuro-symbolic evaluation needs. The ontologies from these benchmarks, together with those from the OWL Reasoner Evaluation (ORE) Competition~\cite{ORE}, can serve as initial datasets for a neuro-symbolic benchmark framework, but do not by themselves address the distinct challenges neuro-symbolic reasoning raises. To the best of our knowledge, no benchmark or evaluation framework has been designed specifically to evaluate and compare neuro-symbolic \emph{description logic} reasoners in general — as distinct from the broader landscape of neuro-symbolic AI benchmarks, some of which (e.g., benchmarks targeting LLM-based multi-step logical reasoning) address adjacent but distinct problems and reasoning targets, and as distinct from noise-robustness specifically, which is the narrower question our own concurrent work addresses (discussed below). Most reasoner evaluations in the DL setting instead rely on ad hoc selections from publicly available ontologies, including but not limited to SNOMED CT\footnote{\url{https://bioportal.bioontology.org/ontologies/SNOMEDCT}}, Gene Ontology (GO)\footnote{\url{https://bioportal.bioontology.org/ontologies/GO}}, and Galen\footnote{\url{https://bioportal.bioontology.org/ontologies/GALEN}}, as well as ontologies from DBpedia~\cite{dbpedia}, YAGO~\cite{yago}, Wikidata~\cite{wikidata}, Claros\footnote{\url{https://www.clarosnet.org}}, NCBO Bioportal\footnote{\url{https://bioportal.bioontology.org/}}, and AgroPortal\footnote{\url{http://agroportal.lirmm.fr/}}. This is a limited set of ontologies that does not cover the full spectrum of possible scenarios.

Table~\ref{tab:nesy} makes this diversity concrete. It shows the use of DL fragments such as $\mathcal{ALC}$ and $\mathcal{EL}^{++}$ and various OWL 2 profiles such as EL and RL~\cite{OWL2profile}, with some works also incorporating RDF and RDFS — underlining a real diversity in supported languages, reasoning tasks, datasets, and evaluation metrics. The table does not aim to be an exhaustive survey; it is included to make the variation in reasoning and evaluation approaches concrete, and to motivate the need for a standardized benchmark supporting fair, consistent comparison. Even superficially similar works can differ substantially: Makni et al.~\cite{noise-tolerant} and Ebrahimi et al.~\cite{pointer_nw} both target RDFS entailment reasoning aiming to replicate deductive reasoning processes, yet adopt different metrics and datasets, making their results difficult to compare directly. Section~\ref{study} takes two of the methods surveyed here — ELEmbeddings~\cite{ELEm} and OWL2Vec*~\cite{owl2vecstar} — and actually runs them under comparable conditions, rather than only citing their reported numbers.

{\small
\begin{longtable}{|p{2.2cm}|p{1.2cm}|p{2.4cm}|p{2.0cm}|p{2.0cm}|p{4.0cm}|}\arrayrulecolor{lightgray}
\hline
\textbf{Paper} & \textbf{Logic} & \textbf{Reasoning Task} & \textbf{Datasets Used} & \textbf{Metrics} & \textbf{Summary of Approaches Used} \\
\hline
ELEm~\cite{ELEm} & $\mathcal{EL}^{++}$ & Subsumption & GO & Hits@n, AUC, Mean Rank & To capture entity relationships, embeddings were created by representing Concepts as n-balls and the relations as translation vectors between the centers of each Concept ball. The embeddings were utilized to predict protein-protein interactions. \\
\hline
EmEL$^{++}$~\cite{Mondal} & $\mathcal{EL}^{++}$ & Subsumption & SNOMED CT, Anatomy, GO, Galen & Hits@n, AUC, Median Rank, 90$^{th}$ percentile rank &
Extended ELEm with relation inclusion and role chains. Also introduced negative samples for training.\\
\hline
EmEL-V~\cite{why_settle_for_one} & $\mathcal{EL}^{++}$ & Subsumption &SNOMED CT,  GO, Galen & Top@n, Median Rank, 90$^{th}$ Percentile Rank & Extended EmEL$^{++}$ to include many-to-many relationships \\
\hline
BoxEL~\cite{faithful_embeddings} & $\mathcal{EL}^{++}$ & Subsumption & Anatomy, GO, Galen & Hits@n, AUC, Mean Rank & To capture entity relationships, mapped concepts as boxes and deals with the limitations of n-ball~\cite{ELEm,Mondal,why_settle_for_one}
based embeddings. \\
\hline

$Box^2EL$~\cite{dualboxembeddings} & $\mathcal{EL}^{++}$ & Subsumption, Role Assertion and Deductive Reasoning & Anatomy, GO, Galen & Hits@n, AUC, Median, Mean Rank & Maps both
concepts and roles as boxes, and
models inter-concept relationships using a bumping mechanism. \\
\hline
{\"{O}}z{\c{c}}ep et al.~\cite{cone_semantics} & $\mathcal{ALC}$ & Concept Membership & NA&NA & Embeds Concepts in the ontology as convex
regions in vector spaces. \\
\hline
E2R~\cite{QuantumEmbedding} & $\mathcal{ALC}$ & Concept Membership & LUBM& Hits@n, Mean Rank, MRR & Aiming to preserve the logical structure, proposed embeddings in the quantum space.\\
\hline
Makni and Hendler~\cite{noise-tolerant} & RDFS & Entailment Reasoning & LUBM and Scientist dataset created from DBpedia&Precision, Recall, and F1 Score & The evaluation focused on assessing noise tolerance by employing an encoder-decoder architecture to translate input RDF graph embeddings into corresponding inference graph embeddings. \\
\hline
Ebrahimi et al.~\cite{Ebrahimi2018ReasoningOR} & RDFS & Query-based Classification & Created from Linked Data Cloud and Data Hub websites& Precision, Recall, and F1 score & Explored the capabilities of end-2-end memory networks. The model's capability for multi-hop reasoning is demonstrated. The use of normalized embeddings support transfer. \\
\hline
Ebrahimi et al.~\cite{pointer_nw} & RDFS and $\mathcal{EL}^{+}$ & Entailment Reasoning & Synthetic Data and LUBM & Exact Matching Accuracy & Utilized pointer networks for learning the sequential application of inference rules used in many deductive reasoning algorithms.\\
\hline
Hohenecker and Lukasiewicz~\cite{patrick} & OWL 2 RL & Entailment Reasoning & Claros, DB-
pedia, UMLS, and Synthetic Data& Accuracy& Developed a deep learning-based model called Recursive Reasoning Networks (RNN).\\
\hline
Eberhart et al.~\cite{emulation} & $\mathcal{EL}^{+}$ & Ontology Completion (concept inclusions and existential restrictions) & Synthetic Data and SNOMED & Precision, Recall, and F1 Score & Showcases completion reasoning behavior using various LSTM neural networks to learn reasoning patterns, employing three distance measures to assess prediction accuracy. \\
\hline
Makni et al.~\cite{Explainable_rdfs_reasoning} & RDFS & Explainable Entailment Reasoning & LUBM and real-world scholarly dataset & Accuracy& Built upon the previous work ~\cite{noise-tolerant} for generating explanations for the derived conclusions by taking the RDF graph and inferred triples as input and the explanations as the target. \\
\hline
Hohenecker and Lukasiewicz ~\cite{DL_for_OR} & RDF & Concept Membership and Relation Prediction & LUBM, UOBM, Claros, DBpedia& F1 Score and Accuracy & Proposed Relational Tensor Network (RTN). Embeddings of the individuals are computed by applying RTNs on the Directed Acyclic Graph representation of the ontology (including the inferences).\\
\hline

Farzana et al.~\cite{query_rewriting} & RDF & Query pruning and complete query rewriting & Created from user search logs from eBay Inc.& Precision, Recall, and F Score, and Query Accuracy & Proposes a Knowledge Graph (KG) enhanced approach for query rewriting in e-commerce, leveraging RDF2Vec entity embeddings, entity types, category information, and entity frequency extracted from a product KG.\\
\hline

OWL2Vec* ~\cite{owl2vecstar} & $\mathcal{SROIQ}$ &Concept membership and concept subsumption &HeLis, FoodOn, GO & MRR and Hits@n & Ontologies are transformed into RDF graphs, and Word2Vec is utilized on the resulting paths. The training dataset comprises three documents: structural, lexical, and a combination of both, enhancing entity interrelation understanding compared to earlier Word2Vec methodologies.\\
\hline

\caption{Overview of Variations in Neuro-Symbolic Reasoning and Evaluation Approaches}

\label{tab:nesy}
\end{longtable}
}

Classifying the works in Table~\ref{tab:nesy} against the six categories above makes the diversity concrete in a second way: \cite{owl2vecstar} converts symbolic input (ontologies, RDF graphs) to vectors (Symbolic Neuro Symbolic). \cite{noise-tolerant,Ebrahimi2018ReasoningOR,pointer_nw,emulation,DL_for_OR,Explainable_rdfs_reasoning} take symbolic reasoning rules as input and compile them during training (Neuro:Symbolic $\rightarrow$ Neuro), integrating symbolic knowledge into neural models. \cite{ELEm,Mondal,why_settle_for_one,faithful_embeddings,cone_semantics,QuantumEmbedding,dualboxembeddings} embed symbolic reasoning inside neural engines (Neuro[Symbolic]). \cite{query_rewriting} falls under the collaborative integration category (Neuro|Symbolic).

\paragraph{A concurrent effort on noise robustness.} One gap the works above share is that none of them tests a system's tolerance to injected noise or inconsistency in a controlled, reproducible way. In separate, concurrent work co-authored with Loesch, Mutharaju, and Celebi~\cite{loesch2024nsorn}, we addressed exactly this gap: NSORN (Neurosymbolic Ontology Reasoning with Noise) introduces three deterministic noise-injection techniques — logical noise (violating asserted \texttt{DisjointClasses} axioms or object-property domain/range constraints), statistical noise (adding low-probability triples predicted by a trained Relational Graph Convolutional Network), and random noise (corrupting the subject or object of existing triples) — and applies them, at five noise levels each, to the OWL2Bench-1 and Family ontologies. It evaluates two of the same systems surveyed in Table~\ref{tab:nesy}, Box2EL~\cite{dualboxembeddings} and OWL2Vec*~\cite{owl2vecstar} (both run through the mOWL library~\cite{mowl}, the same library used in Section~\ref{study}), on class-assertion and object-property-assertion prediction under each noise condition. Its main finding is that logical noise is generally the most damaging of the three, though not uniformly: on OWL2Bench-1, class-assertion MRR for OWL2Vec* drops from 0.070 to 0.043 under 100\% logical noise, while on the Family ontology the same construct barely moves it (0.513 to 0.516), because Family's schema makes some of the noised relations trivially recoverable regardless of the assertions actually corrupted.

NSORN is not itself a contribution of this thesis, and this chapter does not repeat or extend its noise-injection experiments. It is discussed here because it directly narrows what can still be claimed as an open problem in Section~\ref{desiderata}: NSORN already shows that standard MRR/Hits@N metrics \emph{do} shift, in aggregate, under injected inconsistency, for two real systems on two real ontologies at realistic scale (OWL2Bench-1 alone has 60{,}573 axioms). What it explicitly leaves open — by its own account — is noise confined to the TBox (its own techniques target the ABox only, with TBox noise named as future work) and a standardized, cross-system metric suite (also named as future work, given that its own evaluation found the sensitivity of MRR to noise varies enough by ontology and task that a single aggregate number can be misleading). Desideratum~\ref{item:inconsistency} below is written against that more precise boundary rather than against the flat claim that inconsistency-injection is unaddressed.

\section{An Empirical Study: Running Two Surveyed Systems}
\label{study}

The desiderata in Section~\ref{desiderata} could be argued for in the abstract. Instead, this section puts two of the systems from Table~\ref{tab:nesy} to a direct test: ELEmbeddings~\cite{ELEm}, an $\mathcal{EL}^{++}$-restricted embedding method, and OWL2Vec*~\cite{owl2vecstar}, a general-DL ($\mathcal{SROIQ}$) embedding method. The pairing is deliberate. Both are already documented in Table~\ref{tab:nesy}; both have real, maintained implementations in the mOWL library~\cite{mowl} (the same library NSORN uses); and choosing one EL-restricted system and one general-DL system, rather than two systems of the same profile, is what makes the comparison informative for benchmark design specifically — it tests whether a profile boundary that looks like a modelling detail in a paper's related-work section turns into a practical obstacle once you actually hand the system an ontology that crosses it.

Each system was run on three real inputs: two ontologies generated by OWL2Gen-DL (Section~\ref{sec:OWL2Gen_intro}) — one built entirely from OWL2Gen's EL-eligible constructs with the profile set to EL, the other from its full 64-construct catalogue with the profile set to DL — and, for ELEmbeddings, the Anatomy ontology used to evaluate EmEL$^{++}$~\cite{Mondal} and the EL-embedding lineage more broadly, obtained from the EmEL$^{++}$ authors' own repository (Mondal, Bhatia, and Mutharaju, the same research group as this thesis). All code, generation parameters, and raw logs for this study are collected in a companion repository (see the note at the end of this section); nothing reported below is fabricated or interpolated from partial runs.

\subsection{Setup}

The two OWL2Gen-DL ontologies were generated with \emph{Guarantee consistency} enabled, seed 2026, and default hierarchy and nesting parameters. The EL-profile request asked for 80 occurrences of each of OWL2Gen's 30 EL-eligible constructs against the OWL~2 EL profile; the DL-profile request asked for 25 occurrences of each of all 64 supported constructs against the OWL~2 DL profile. The Anatomy ontology, reused as-is from EmEL$^{++}$'s own experimental data, comprises 266{,}701 axioms over 103{,}721 classes and 191 object properties, merging UBERON, FMA, and species-specific anatomy ontologies under a single EL-normal-form export. Loading it required two small, mechanical repairs to the distributed file: it ships as a bare sequence of per-axiom lines without the \texttt{Ontology(\ldots)}/\texttt{Prefix(\ldots)} envelope OWL~API parsers require (its original codebase uses its own line-based parser rather than the OWL~API), and four axioms use a degenerate empty \texttt{ObjectPropertyChain()} that no OWL~2 parser accepts. Wrapping the file in a standard envelope and dropping those four axioms — out of 266{,}701 — was sufficient to load it correctly; neither change alters the ontology's actual content.

For each input, both systems were trained under a single held-out evaluation of asserted, atomic \texttt{SubClassOf} axioms (i.e., $A\sqsubseteq B$ with $A$ and $B$ both named classes): a sample of these axioms was removed before training and then ranked for recovery under the standard filtered protocol, following the same methodology as the pilot study in the previous version of this chapter. ELEmbeddings was scored with its own trained geometry — the $\mathcal{EL}^{++}$ ball-inclusion score $\lVert c_A-c_B\rVert+r_A-r_B$ between class centers $c$ and radii $r$ — while OWL2Vec* was scored by cosine similarity between its random-walk-derived class embeddings, matching how each method is used for subsumption prediction in its own literature. Anatomy was trained for 3{,}000 epochs (a full pass over 266{,}701 axioms, including 109{,}605 atomic \texttt{SubClassOf} facts, took under a quarter of a second per epoch); the two much smaller OWL2Gen ontologies were trained for 2{,}000 epochs (ELEmbeddings) or 20 word2vec epochs over 20 random walks per node (OWL2Vec*). One run, one seed, per input, matching this chapter's earlier pilot's own stated scope: this is a small, targeted study, not a full benchmark.

\subsection{Finding 1: Requesting a Profile Does Not Guarantee It}
\label{sec:study-profile}

The first result did not come from either neuro-symbolic system at all — it came from OWL2Gen-DL's own verification step. The EL-profile ontology, built exclusively from constructs OWL2Gen itself labels EL-eligible and generated with the profile explicitly set to EL, was correctly reported consistent by Openllet, but its OWL~2 profile classification came back \texttt{el: false, ql: false, rl: false, dl: false} — outside every one of the four OWL~2 profiles, including plain OWL~2 DL, with the OWL~API's own profile checker reporting 38 violations, all of the form "use of non-simple property in \texttt{ObjectHasSelf} restriction." Tracing the cause required going past a first, plausible-looking guess: the two properties involved (\texttt{objectProperty5} and \texttt{objectProperty9}) were never themselves declared transitive, and OWL2Gen-DL's local consistency guard does correctly prevent \texttt{ObjectHasSelf} from picking a \emph{directly} transitive property. What the guard missed is indirect non-simplicity: the same request separately declared \texttt{EquivalentObjectProperties(objectProperty2, objectProperty5)} and \texttt{EquivalentObjectProperties(objectProperty13, objectProperty9)}, and \texttt{objectProperty2} and \texttt{objectProperty13} were both independently declared transitive elsewhere in the same ontology. OWL~2 non-simplicity propagates through declared equivalence — two properties declared equivalent must denote the same relation, so one cannot be simple while the other is transitive — but OWL2Gen-DL's guard tracked this propagation only through declared inverses, not through \texttt{EquivalentObjectProperties} groups, so \texttt{objectProperty5} and \texttt{objectProperty9} still looked simple to the \texttt{ObjectHasSelf} generator at the point they were picked. This is now fixed directly in the tool: \texttt{ConstraintTracker}'s simple/non-simple tracking propagates through equivalence groups in both directions (marking and eligibility checks), and \texttt{EquivalentObjectProperties} generation itself now refuses to merge a group that would combine an already-non-simple property with an already-simple-required one, closing the gap at the point the equivalence axiom would have been added rather than only detecting the resulting ontology as broken afterward. The fix was verified at two levels, not only by inspection: a standalone reproduction against the tracker logic directly confirmed that, before the fix, a property merged by equivalence into an already-transitive property's group was still reported eligible for \texttt{ObjectHasSelf}, and that after the fix it is correctly rejected, with the reverse-order case (equivalence-merging into an already-simple-required property) also now correctly refused; and, end-to-end, reissuing the identical EL-profile request against the rebuilt tool now yields an ontology the OWL~API classifies as \texttt{el: true, dl: true} with \texttt{profileGuaranteeSatisfied: true} — a complete reversal of the original \texttt{el: false, ql: false, rl: false, dl: false} result, on the same request.

This matters beyond being a bug report. Every individual axiom in the ontology was drawn from a construct OWL2Gen labels EL-eligible, the profile was explicitly requested, and the result was still not just non-EL but non-DL. If a purpose-built generator with an explicit profile-guard mechanism can produce this outcome, an ad hoc selection from an existing ontology — the norm across Table~\ref{tab:nesy} — offers no guarantee at all that a dataset advertised as "EL" or "DL" actually is. Section~\ref{desiderata}'s desideratum on diverse, profile-varying scenarios is written with this specific failure mode in mind: profile conformance has to be verified, not assumed from how the request was phrased.

The DL-profile ontology (1{,}704 axioms after the completeness pass, from a 1{,}600-axiom request) did not have this problem — it was correctly classified as OWL~2 DL and not EL/QL/RL — but its consistency could not be established within the 30-second budget (Openllet reported \texttt{NOT\_VERIFIED\_TIMEOUT} after 30{,}009\,ms). This is the same expressivity-driven reasoning cost documented for classical reasoners in Section~\ref{sec:OWL2Gen_intro}, here surfacing again as a precondition for the neuro-symbolic experiments below: neither ELEmbeddings nor OWL2Vec* requires a prior consistency verdict to train, so both experiments proceeded regardless, but a benchmark that wanted to report "trained on a consistent DL ontology" as a controlled variable would not have been able to certify that within a normal time budget here.

\subsection{Finding 2: The Same Input, Two Different Failure Modes}
\label{sec:study-representation}

The clearest result of this study came from feeding the DL-profile ontology to ELEmbeddings. The training pipeline did not silently degrade — it crashed outright, before producing a single embedding:

\begin{quote}
\texttt{de.tudresden.inf.lat.jcel.owlapi.translator.TranslationException: This axiom is not supported: 'DatatypeDefinition(<http://owl2gendl.com/generated\#datatype39> DataIntersectionOf(rdfs:Literal xsd:string))'}
\end{quote}

The offending axiom is a single \texttt{DatatypeDefinition}, one of OWL2Gen's deliberately non-EL constructs, and a perfectly legal OWL~2 DL axiom. mOWL's $\mathcal{EL}^{++}$ normalizer (built on the jcel reasoner's translator) has no representation for it and raises an unhandled exception rather than skipping it, which means an EL-restricted embedding method cannot be pointed at an arbitrary DL-profile ontology at all without a preprocessing step to strip or rewrite whatever falls outside its supported fragment — a step no benchmark in Table~\ref{tab:nesy} reports performing or characterizing.

OWL2Vec* trained on the identical file without error, since its graph-projection step has no notion of an EL normal form to conform to. But it did not fully succeed either, just more quietly: of the 285 classes in the ontology, 201 received a walk-derived embedding, and none of the 5 held-out \texttt{SubClassOf} axioms' subject classes were among them, so zero held-out queries could be ranked at all. The gap is a coverage artifact of short random walks on a small, sparsely-connected graph rather than a crash, but the practical consequence for a benchmark is the same: a result table that only reports "trained successfully" would treat this run as usable, when in fact no subsumption could be evaluated from it.

Two systems, the same input, two different failure modes — one loud and total, one quiet and partial. Neither is visible from the systems' own published evaluations, which never expose them to an ontology outside their respective comfort zones. This is the concrete version of Section~\ref{desiderata}'s desideratum on input representation: profile and format compatibility is not a preprocessing footnote, it is a variable that determines whether an experiment produces any usable result at all.

\subsection{Finding 3: What the Numbers Show Where They Exist}

Table~\ref{tab:nesy-validation-study} reports every filtered-ranking result this study could actually compute. Where a cell is dashed, the corresponding run either crashed (Finding 2) or produced zero evaluable queries (Finding 2, OWL2Vec* on the DL-profile ontology); those cells are not zeroes, they are the absence of a number, and are reported as such rather than folded into an average.

\begin{table}[htbp]
\centering
\caption{Filtered-ranking recoverability of held-out \texttt{SubClassOf} axioms, using each method's own scoring function (ball-inclusion distance for ELEmbeddings, cosine similarity for OWL2Vec*). A dash means the run did not produce a usable result (see Section~\ref{sec:study-representation}), not a score of zero.}
\label{tab:nesy-validation-study}
\begin{tabular}{llrrrrr}
\toprule
Method & Input & Held-out & MRR & Hits@1 & Hits@3 & Hits@10 \\
\midrule
ELEmbeddings & Anatomy (266{,}701 ax.)      & 227 & 0.0090 & 0.000 & 0.004 & 0.013 \\
ELEmbeddings & OWL2Gen EL-profile (2{,}213 ax.) & 16  & 0.0081 & 0.000 & 0.000 & 0.000 \\
ELEmbeddings & OWL2Gen DL-profile (1{,}704 ax.) & --  & --     & --    & --    & --    \\
\midrule
OWL2Vec* & OWL2Gen EL-profile (2{,}213 ax.) & 16 & 0.0358 & 0.000 & 0.000 & 0.125 \\
OWL2Vec* & OWL2Gen DL-profile (1{,}704 ax.) & 0  & --     & --    & --    & --    \\
\bottomrule
\end{tabular}
\end{table}

Three observations follow, none of which was assumed going in. First, where both methods could be run on the same input (the EL-profile ontology), OWL2Vec* substantially outperforms ELEmbeddings on Hits@10 (0.125 vs.\ 0.000) despite ELEmbeddings being the profile-matched, purpose-built method for exactly this kind of axiom. This is not evidence that OWL2Vec* is "better" in any general sense — the two methods were trained under very different budgets (2{,}000 gradient-descent epochs against 16 held-out queries for ELEmbeddings vs.\ 20 word2vec epochs over 20 walks per node for OWL2Vec*, on an ontology with only 280 classes and 80 atomic subsumption facts total) — but it is a direct illustration of why Section~\ref{desiderata}'s desideratum on metrics argues against reporting a single method's numbers as representative of "what neuro-symbolic embedding can do" for a given profile: architecture and training budget move the number as much as, or more than, profile match does.

Second, scale does not rescue recoverability by itself. ELEmbeddings on Anatomy — 121 times more axioms than the OWL2Gen EL-profile ontology, and 3{,}000 training epochs instead of 2{,}000 — reaches Hits@10 of only 0.013, worse in absolute terms than the small ontology's 0.000-but-comparable MRR (0.0090 vs.\ 0.0081). More data and more training epochs did not translate into materially better recoverability at this training budget; a benchmark that used raw axiom count as a proxy for "how hard this dataset is for embedding methods" would not have predicted this.

Third, and consistent with the analogous finding in the earlier version of this pilot: neither method comes remotely close to what a symbolic reasoner does on the same input. Openllet verified the EL-profile ontology's consistency in 31\,ms and would compute its full deductive closure exactly, with no training at all, from the same 2{,}213 axioms neither embedding method could recover more than 12.5\% of ten guesses on. The point is not that embedding methods are bad at description logic — it is that a benchmark reporting only embedding-side metrics, without a symbolic reference point run on the identical input, cannot communicate how much recoverability is actually being left on the table.

\subsection{Summary and Scope}

This study is deliberately small: two systems, three inputs, one seed each, chosen because they are already documented in Table~\ref{tab:nesy} and have real, runnable implementations rather than because they are the strongest possible baselines. Its purpose is not to rank ELEmbeddings against OWL2Vec*, but to show that the difficulties Sections~\ref{desiderata} and~\ref{benchmark} argue for are not hypothetical: within a few hours of actually running two existing, published, well-cited systems side by side, this study encountered a profile-verification gap in ontology generation itself, a hard crash and a silent coverage gap from the same input on two different systems, and metrics that move more with architecture and training budget than with the profile-match a reader might expect to be the deciding factor. All generation parameters, raw ontology files (including the repaired Anatomy export), training scripts, and logs for this study are collected in a standalone companion repository, separate from the main OWL2Gen-DL codebase, so the runs can be reproduced or extended independently.

\section{Desiderata for Benchmarking Neuro-Symbolic Reasoners}
\label{desiderata}

Creating an effective benchmark demands careful consideration of principles such as simplicity for accessibility, portability for impartial assessment across approaches, scalability to accommodate diverse system sizes, and relevance to practical challenges~\cite{jimgray}. The evaluation of neuro-symbolic reasoners raises additional, distinctive challenges that the field's novelty means current state-of-the-art work does not yet address systematically. We outline six issues that a fair neuro-symbolic reasoning benchmark should prioritize, and, where Section~\ref{study} or~\cite{loesch2024nsorn} encountered the issue directly rather than only motivating it in the abstract, we say so.

\begin{enumerate}
    \item \textbf{Diverse benchmark scenarios, with verified profiles.} Neuro-symbolic reasoning aims to make sense of complex application domains, so a benchmark should reflect that complexity with ontologies that vary in size (ABox and TBox), profile, and axiom types. Because different neuro-symbolic reasoners work in different ways, benchmark scenarios should include both specific reasoning tasks — enabling micro-benchmarking — and generic information needs that can be implemented differently across architectures, ensuring comparability across heterogeneous systems. Section~\ref{sec:study-profile} shows this cannot stop at requesting a profile: an ontology built entirely from constructs a generator itself labels EL-eligible, with the profile explicitly requested, still verified as outside every OWL~2 profile once its axioms interacted with each other. A benchmark's scenario-generation step therefore needs to verify the profile of what it actually produced, not the profile of what it asked for.

    \item \label{item:inconsistency} \textbf{Introducing controlled inconsistencies, with metric sensitivity reported.} A benchmark should be able to generate controlled inconsistencies deterministically — for example, asserting that an instance belongs to two TBox-disjoint classes, or introducing near-duplicate entity names (e.g., "John" vs.\ "Jonh") to test semantic-inconsistency handling. This evaluates a system's ability to detect and resolve contradictions in a way that mirrors real-world conflicting information. Our own concurrent work, NSORN~\cite{loesch2024nsorn}, is a first step in exactly this direction: it introduces logical, statistical, and random ABox noise deterministically and at controlled levels, and evaluates two neuro-symbolic reasoners against it on two real ontologies. What it also shows, by its own account, is that the sensitivity of standard metrics to that noise is inconsistent — logical noise degrades OWL2Vec*'s class-assertion MRR by nearly half on OWL2Bench-1 but leaves it almost unchanged on Family, because Family's own schema makes some corrupted relations trivially recoverable regardless. A benchmark that introduces controlled inconsistencies therefore needs to report, as a first-class result and not an aside, whether and how strongly its chosen metrics actually respond to the inconsistency introduced — since NSORN's own data shows that response is not a constant, and its own stated future work leaves TBox-level noise injection and a standardized cross-system metric suite open.

    \item \textbf{Input representation, checked for compatibility, not just diversity.} A central challenge in neuro-symbolic reasoning is representing ontological information as embeddings. Traditional techniques such as TransE~\cite{Bordes_EL+}, though powerful, were not designed to retain logical information, which typically forces a separation between ABox (assertional) and TBox (terminological) knowledge during embedding. Engineering a representation that captures logical relationships is a prerequisite for a well-founded neuro-symbolic reasoner. Section~\ref{sec:study-representation} shows what happens when this is left unchecked: an EL-restricted embedding method did not degrade gracefully on a legal OWL~2 DL axiom outside its supported fragment, it crashed outright, while a graph-projection-based method tolerated the same input but silently lost embedding coverage for enough entities that no held-out query could be evaluated at all. A benchmark should generally avoid imposing a specific input representation unless the explicit goal is testing a particular neuro-symbolic property, but it must report, for every system under test, whether the representation step itself succeeded on a given input and what was lost or dropped if it did not — a fact neither crash nor silent coverage loss will surface on its own.

    \item \textbf{Assessment of deductive capabilities, against a symbolic reference point.} Building a new generation of reasoners that harness both neural networks and logical reasoning requires an equitable assessment of existing state-of-the-art solutions. Current approaches show promising but limited deductive capabilities, and it is important to understand how far these capabilities can be pushed. Section~\ref{study}'s third finding puts a number on this gap directly: Openllet computed the exact deductive closure of the EL-profile ontology in 31\,ms with no training at all, on the same input where neither embedding method exceeded 12.5\% Hits@10. A comprehensive evaluation should also include purely neural approaches, such as large language models, as a baseline: even where such approaches do not excel on deduction, including them clarifies the comparative landscape and any synergies between purely neural and neuro-symbolic methods.

    \item \textbf{Success metrics that account for architecture and training budget, not just paradigm.} Neuro-symbolic evaluation needs metrics that combine the logical principles of soundness and completeness with statistical measures of performance, rather than relying on either alone. Section~\ref{study}'s comparison of ELEmbeddings and OWL2Vec* on the identical ontology shows why a single method's Hits@n cannot stand in for "what neuro-symbolic embedding can achieve" on a given profile: the profile-matched, purpose-built method performed worse than the general-purpose one, a difference plausibly driven more by training budget and architecture than by profile alignment. An ideal neuro-symbolic reasoner would support the most expressive logic, transfer across domains, generate all and only correct inferences in a single run while remaining sound, and explain its conclusions at scale with good performance. A benchmark should evaluate the degree to which systems approach these properties while varying architecture and training budget deliberately, since Section~\ref{study} shows both are large enough to change the ranking between systems on their own.

    \item \textbf{Adaptability.} Given how quickly the field is evolving, a benchmark needs to accommodate new tasks and evolving requirements without becoming stale, remaining flexible enough to reflect domain-specific nuances as they emerge. The EL-embedding lineage surveyed in Table~\ref{tab:nesy} alone moved from ELEm~\cite{ELEm} to EmEL$^{++}$~\cite{Mondal} to BoxEL~\cite{faithful_embeddings} to Box$^2$EL~\cite{dualboxembeddings} within a few years, each changing the geometry and the claims made about it; a benchmark that hard-codes assumptions about any one of these will be stale before it is finished. This keeps the benchmark useful as a long-term point of comparison rather than a snapshot of one moment in the field.
\end{enumerate}

\section{Methodology for Designing the Benchmark}
\label{benchmark}

This section proposes one possible methodology for designing a benchmark that addresses the desiderata of Section~\ref{desiderata}, revised where Section~\ref{study} changed what the methodology needs to guard against. The methodology comprises six steps.

\begin{enumerate}
\item \textit{Dataset curation for diverse scenarios, with profile verification as a build step.} Curating datasets that cover a wide range of benchmarking scenarios starts with a study of existing neuro-symbolic DL reasoners, beginning with basic ontology profiles such as RDFS and progressing to more expressive ones such as OWL 2 EL and OWL 2 DL. This includes evaluating existing datasets across various models and systems, comparing results against traditional reasoning systems such as Konclude~\cite{konclude}, and identifying performance variation. Section~\ref{sec:study-profile} shows that requesting a profile is not the same as producing one: profile conformance of every generated (or curated) scenario needs to be checked by an actual OWL~2 profile classifier as part of dataset curation, not assumed from the generation request or the source ontology's label. Diversity should also be measured structurally, not only by size and profile: an earlier version of this pilot showed that two evaluation sets with identical size and identical asserted/inferred composition can still differ substantially in average difficulty depending on the connectivity of their target entities, so dataset curation should stratify evaluation splits by structural connectivity (e.g., node in-degree) rather than sample them uniformly at random.

\item \textit{Introduction of controlled inconsistencies, building on NSORN.} After curating the dataset, controlled inconsistencies should be introduced to assess system robustness. NSORN~\cite{loesch2024nsorn} already provides deterministic logical, statistical, and random ABox noise-injection techniques applicable at this step; a full benchmark should adopt and extend them to the TBox, which NSORN's own account leaves as future work. Equally important, and easy to overlook, is verifying that the benchmark's own reported metrics can actually detect the inconsistency once introduced: NSORN's data already shows this sensitivity varies by ontology and task rather than holding uniformly, so a benchmark that introduces controlled inconsistencies should report, as a first-class result, how strongly its chosen metrics respond to them on each dataset, not only whether the injection mechanism ran.

\item \textit{Appropriate input formats, with compatibility as a reported outcome.} Given how varied neuro-symbolic reasoners' input requirements are, an exhaustive list of required formats is impractical. A more tractable solution is a tool supporting ontology generation in common user-required formats such as RDF/XML\footnote{\url{https://www.w3.org/TR/rdf-syntax-grammar/}} or Manchester OWL Syntax\footnote{\url{https://www.w3.org/2007/OWL/wiki/ManchesterSyntax}}, together with configurable dataset-split generation to meet diverse testing needs. Section~\ref{sec:study-representation} shows this is not only a formatting concern: an EL-restricted system can fail outright, and a general-DL system can silently lose entity coverage, on the identical legally-formatted input. The benchmark harness should therefore run a compatibility pre-check for each system under test — attempting the representation step and recording exactly what failed or was dropped — before any downstream metric is computed, rather than only reporting metrics for the runs that happened to succeed.

\item \textit{Evaluation of deductive capabilities.} Because traditional reasoners often struggle with inconsistent ontologies, generating inconsistencies should remain a separate step, with deductive evaluation conducted first on consistent ontologies. This step should assess the neural component's learning capability, repair ability, and scalability — handling ontological complexity, scaling behavior, and overall performance relative to traditional reasoning systems, reported alongside the symbolic reasoner's own result on the identical input as Section~\ref{study} did. Repeating this evaluation after controlled inconsistencies have been introduced highlights the specific benefits the neural component contributes under imperfect information, giving a more complete picture of system capability than either evaluation alone.

\item \textit{Metric design, controlling for architecture and training budget.} Standard learning metrics such as accuracy, precision, and F1 score give only an overall picture of system efficacy. Section~\ref{study} shows that a single method's numbers on a given profile can be driven as much by architecture and training budget as by anything specific to the profile itself, so a benchmark's metric-reporting protocol should record and vary these factors deliberately — for example by running more than one architecture per profile, as Section~\ref{study} did — rather than let them vary incidentally between runs. A more thorough analysis also needs metrics that separate out which ontological constructs and reasoning patterns a system handles well versus poorly, and that capture adaptability to diverse scenarios in addition to raw deductive performance.

\item \textit{Ensuring adaptability.} Finally, the benchmark's dataset and task set should be updated regularly as new systems, reasoning tasks, and input formats emerge, so that it remains relevant and usable over time.
\end{enumerate}

This methodology is proposed at the level of principle rather than as a complete, implemented benchmark — building and validating the full pipeline described above is out of scope for this thesis. Section~\ref{study} is a targeted demonstration that the gaps this methodology is designed to close are real and encountered quickly in practice, not a substitute for the full benchmark itself.

\section{Conclusion and Future Work}
\label{conclusion}

This chapter argued that benchmarking neuro-symbolic description logic reasoners requires evaluation principles fundamentally different from those developed for classical symbolic reasoners, and it grounded that argument in a direct test rather than leaving it at the level of survey and assertion. Section~\ref{study} took two systems already documented in Table~\ref{tab:nesy} — an EL-restricted method and a general-DL method, both with real, runnable implementations — and ran them on OWL2Gen-DL output and on a real ontology reused from the EL-embedding literature. Within that small study, requesting an OWL~2 EL profile did not guarantee producing one; the same DL-profile ontology crashed one system outright and silently starved the other of usable coverage; and the two systems' recoverability numbers on the one input they shared moved more with architecture and training budget than with which method was profile-matched to the task. None of this was assumed going in, and none of it is visible from the systems' own published evaluations, which do not expose them to inputs outside their comfort zones. Sections~\ref{desiderata} and~\ref{benchmark} state the resulting desiderata and a possible benchmark-design methodology directly against these findings, and against the concurrent noise-robustness work in~\cite{loesch2024nsorn}, which is related to this chapter's argument but is not itself part of this thesis's contribution.

The contributions of this thesis lay a foundation for future work along several lines suggested directly by Section~\ref{study} as well as by the broader methodology. First, Section~\ref{study}'s findings should be stress-tested beyond three inputs and one seed each: repeating the profile-verification check across more OWL2Gen-DL configurations, and matching ELEmbeddings' and OWL2Vec*'s training budgets more closely, would clarify how much of the observed gap between them is architectural versus incidental. Second, the OWL2Gen-DL construct-compatibility gap identified in Section~\ref{sec:study-profile} — non-simplicity not propagating through declared \texttt{EquivalentObjectProperties} groups — has since been corrected directly in the tool's shared \texttt{ConstraintTracker} logic and verified against the specific scenario that exposed it. Because every OWL2Gen-DL construct subject to the same simple-property restriction (the cardinality restrictions, \texttt{FunctionalObjectProperty}, \texttt{InverseFunctionalProperty}, \texttt{IrreflexiveProperty}, \texttt{AsymmetricProperty}, and \texttt{ObjectPropertyDisjointWith}, in addition to \texttt{ObjectHasSelf}) already routes through that same shared tracking, the fix covers all of them, not only the construct through which the gap was first observed. Third, the methodology of Section~\ref{benchmark} — dataset curation with an explicit profile-verification step, controlled inconsistency injection extended to the TBox and reported for metric sensitivity, and a compatibility pre-check reported ahead of any downstream metric — should be scaled from this study into the full benchmark framework it was designed to validate.
